La hipótesis de De Broglie, formulada en 1924 \cite{debroglie1924}, establece que las partículas materiales poseen propiedades ondulatorias, asociando a cada partícula una longitud de onda inversamente proporcional a su cantidad de movimiento. Este postulado fue verificado experimentalmente en 1927 mediante la observación de la difracción de electrones al atravesar cristales \cite{davisson1927}.

La difracción es un fenómeno característico de las ondas, en el cual la propagación de una onda se ve alterada al encontrarse con obstáculos \cite{thomson1928}, generando nuevos frentes de onda que pueden interferir entre sí. En el caso de los electrones, este comportamiento se evidenció en experimentos que mostraron patrones de interferencia similares a los observados con ondas luminosas (véase la figura \cref{fig:bragg}).

\begin{figure}[htbp!]
  \begin{center}
    \includegraphics[width=0.85\linewidth]{Images/diffracion.png}
    \caption{Refracción en ondas.}
    \label{fig:bragg}
  \end{center}
\end{figure}

La relación de De Broglie permite calcular la longitud de onda de una partícula a partir de su cantidad de movimiento $p$, según la expresión $\lambda = h / p$, donde $h$ es la constante de Planck. En términos de la energía $E$, esta relación puede reescribirse como $\lambda = hc / E$.

Para un electrón acelerado por una diferencia de potencial $U$, su energía total incluye tanto su energía en reposo como la energía adquirida del campo eléctrico. Al aplicar la relación relativista entre energía y momento, dada por \cref{eq:Energy},

\begin{equation}\label{eq:Energy}
    (pc)^2 = (mc^2 + eU_a)^2 - (mc^2)
\end{equation}

se obtiene una expresión modificada para la longitud de onda:

\begin{equation}\label{eq:lambda}
  \lambda = \frac{h c}{\sqrt{(e U)^2 + 2 e U m c^2}},
\end{equation}

la cual muestra explícitamente la dependencia de la longitud de onda con la energía adquirida por el electrón.

Para que sea posible la difracción de los electrones, la longitud de onda \cref{eq:lambda} debe ser comparable con la distancia interplanar $d$ visto en la figura \cref{fig:graphite}.

\begin{figure}[htbp!]
  \begin{center}
    \includegraphics[width=0.85\linewidth]{Images/graphite.png}
    \caption{Red cristalina de grafito.}
    \label{fig:graphite}
  \end{center}
\end{figure}

Con la longitud de onda dada por \cref{eq:lambda}, se puede aplicar la ley de difracción de Bragg dado que el electrón se comporta como se observa en la figura \cref{fig:diffraction}

\begin{figure}[htbp!]
  \begin{center}
    \includegraphics[width=0.85\linewidth]{Images/diffracion.png}
    \caption{Refracción en ondas al cruzar por dos rendijas.}
    \label{fig:diffraction}
  \end{center}
\end{figure}

La ley de bragg \cite{bragg1913} para el comportamiento de la figura \cref{fig:bragg} se modela en la ecuación \cref{eq:Bragg}.

\begin{equation}\label{eq:Bragg}
    2d\sin{\theta} = n\lambda,
\end{equation}

Donde $d$ es la distancia reticular interplanar del grafito y $\theta$ es el ángulo de refracción. Considerando ángulos pequeños, se puede aproximar $\sin{2\theta} = 2\sin{\theta} = \frac{r}{R}$, y al despejar en \cref{eq:Bragg}, se obtiene una ecuación lineal de la forma $a\lambda + b$:

\begin{equation}\label{eq:linear_r}
    r = \frac{2R}{d} n \lambda.
\end{equation}

Para determinar la distancia interplanar del material, se puede aplicar la ecuación \cref{eq:linear_r}, lo que lleva a la expresión para el parámetro esférico en \cref{eq:parameter-spherical} y para el parámetro lineal en \cref{eq:parameter-linear}:

\begin{equation}\label{eq:parameter-spherical}
  d = \frac{2 R}{a},
\end{equation}

\begin{equation}\label{eq:parameter-linear}
  d = \frac{L}{a}.
\end{equation}
